\worldchapter{Charaktere}{character}
\label{ch:charaktere}

\begin{multicols*}{2}
\section{Was einen Charakter ausmacht}\label{sec:was-einen-charakter-ausmacht2}

Ein Charakter wird auf seinem Charakterbogen in Zahlen und anderen Angaben festgehalten.
Dabei besteht der Charakter aus den folgenden Angaben:

\begin{itemize}[leftmargin=*]
    \item einem Namen
    \item dem Geburtsdatum, aus dem sich auch das Jahrhundert des Spiels ergibt
    \item drei Start-Prägungen: \textbf{Abstammung}\index{Abstammung}, \textbf{Weg}\index{Weg}, \textbf{Bindung}\index{Bindung}
    \item einem Konzeptsatz
    \item sechs Attributen
    \item dreizehn Fertigkeiten
    \item optionalen Zauberschulen
    \item abgeleiteten Werten
    \item den Leisten \textbf{Wunden}\index{Wunden}, \textbf{Bürde}\index{Bürde}, \textbf{Omen}\index{Omen} sowie gegebenenfalls \textbf{Arkana}\index{Arkana}, \textbf{Gunst}\index{Gunst} und \textbf{Verderbnis}\index{Verderbnis}
\end{itemize}

\section{Einen Charakter erstellen}\label{sec:einen-charakter-erstellen}

Um einen Charakter zu erstellen, folge den zehn Schritten in der angegebenen Reihenfolge.
Jeder Schritt baut auf den vorherigen auf und führt dich durch die Definition deines Charakters.

\subsection{Schritt 1: Konzept}\label{subsec:schritt-1:-konzept}

Der Konzeptsatz dient dazu, dir eine Vorstellung von deinem Charakter zu geben.
Er besteht aus einem kurzen Satz oder einer Phrase, die beschreibt, wer dein Charakter ist.
Das Konzept hilft dir dabei, Entscheidungen während des Spiels im Einklang mit der Persönlichkeit deines Charakters zu treffen.
Es muss nicht mechanisch präzise sein, sondern soll als Inspiration dienen.

Beispiele:
\begin{itemize}[leftmargin=*]
\item \textit{Ein abtrünniger Inquisitor auf der Suche nach Vergebung}
\item \textit{Eine Kräuterkundige, die zu viel über verbotene Rituale weiß}
\item \textit{Ein ehemaliger Soldat, der seine Kameraden im Stich ließ}
\item \textit{Eine Adlige, die ihr Erbe für verbotenes Wissen opferte}
\end{itemize}

\subsection{Schritt 2: Geburtsdatum}\label{subsec:schritt-2:-geburtsdatum}

Die Kampagne spielt in einem Jahrhundert von \textbf{1} bis \textbf{10}, vorgegeben von der Spielleiterin.
Daraus ergeben sich \textbf{Glaubensniveau}\index{Glaubensniveau} und \textbf{Magieniveau}\index{Magieniveau}.

Wähle ein Geburtsdatum innerhalb des vorgegebenen Jahrhunderts.

\begin{center}
\begin{tabular}{ccc}
\toprule
\textbf{Jahrhundert} & \textbf{Glaube} & \textbf{Magie} \\
\midrule
1 & 1 & 1 \\
2 & 1 & 2 \\
3 & 1 & 3 \\
4 & 1 & 4 \\
5 & 1 & 5 \\
6 & 2 & 4 \\
7 & 3 & 3 \\
8 & 4 & 2 \\
9 & 5 & 1 \\
10 & 6 & 0 \\
\bottomrule
\end{tabular}
\end{center}

\subsection{Schritt 3: Attribute}\label{subsec:schritt-3:-attribute}

Die Attribute werden von 0 bis 3 vergeben:

\begin{description}[style=nextline,leftmargin=0pt]
\item[\textbf{Geist}] Fähigkeit zur Analyse komplexer Zusammenhänge, Erinnerungsvermögen an wichtige Details und Verständnis okkulter Logik sowie magischer Prinzipien
\item[\textbf{Wille}] Mentale Disziplin und Selbstbeherrschung, Durchhaltevermögen in schwierigen Situationen und die Kraft der Überzeugung im Glauben
\item[\textbf{Instinkt}] Scharfe Wahrnehmung der Umgebung, schnelle Reflexe in Gefahrensituationen und ein natürlicher Sinn für drohende Bedrohungen
\item[\textbf{Geschick}] Präzision bei feinmotorischen Aufgaben, körperliche Geschwindigkeit und Beweglichkeit sowie Kontrolle über den eigenen Körper
\item[\textbf{Körper}] Rohe physische Kraft, körperliche Zähigkeit und Ausdauer sowie natürlicher Widerstand gegen Schaden und Erschöpfung
\item[\textbf{Erscheinung}] Persönliche Ausstrahlung und Präsenz, ausgeübter sozialer Druck auf andere sowie allgemeine Wirkung in gesellschaftlichen Situationen
\end{description}

Setze:
\begin{itemize}[leftmargin=*]
\item ein Attribut auf 3
\item zwei Attribute auf 2
\item zwei Attribute auf 1
\item ein Attribut auf 0
\end{itemize}

Dies sind deine finalen Attributswerte.
Prägungen können später im Spiel zusätzliche Boni auf einzelne Attribute gewähren.

\subsection{Schritt 4: Fertigkeiten}\label{subsec:schritt-4:-fertigkeiten}

Die Fertigkeiten lauten:\\[0.3cm]

\begin{multicols*}{2}
\begin{itemize}[leftmargin=*]
\item Athletik
\item Kampf
\item Schießen
\item Heimlichkeit
\item Handwerk
\item Heilen
\item Wahrnehmung
\item Mythen
\item Einfluss
\item Standhaftigkeit
\item Zauberei
\item Ritus
\item Magiekunde
\end{itemize}
\end{multicols*}

Setze:
\begin{itemize}[leftmargin=*]
\item eine Fertigkeit auf Rang 3
\item drei Fertigkeiten auf Rang 2
\item funf Fertigkeiten auf Rang 1
\item alle übrigen auf 0
\end{itemize}

\subsection{Schritt 5: Prägungen}\label{subsec:schritt-5:-pragungen}
Wähle:
\begin{itemize}[leftmargin=*]
\item eine \textbf{Abstammung}
\item ein \textbf{Weg}
\item eine \textbf{Bindung}
\end{itemize}

Neue Charaktere starten \textbf{ohne Mal}.
Male entstehen erst im Spiel, etwa nach einem Wendepunkt oder überlebtem Sterben.

Mögliche \textbf{Prägungen}\index{Prägung} stehen gesammelt in \cref{app:marks}.

\subsection{Schritt 6: Schuld oder Eid}\label{subsec:schritt-6:-schuld-oder-eid}

Wenn dein Charakter einen Eid geschworen oder eine Schuld zu begleichen hat, führe sie hier auf.
Wähle einen konkreten Eid oder eine Schuld, die mit einer Bindung oder einem Weg verknüpft ist.
Der Eid kann dabei frei formuliert werden, es gibt keine Liste.

\begin{itemize}[leftmargin=*]
\item Wenn der Eid oder die Schuld das Spiel signifikant erschwert: +1 Omen
\item Wenn der Eid oder die Schuld deinem Charakter dienlich oder nützlich ist: +1 Bürde
\end{itemize}

Beispiele für erschwerende Eide und Schulden (+1 Omen):
\begin{itemize}[leftmargin=*]
\item \textit{Eid, niemals eine bestimmte Stadt zu betreten}
\item \textit{Schuld gegenüber einem verfeindeten Adeligen}
\item \textit{Eid, eine verfluchte Reliquie an einen fernen Ort zu bringen}
\end{itemize}

Beispiele für nützliche Eide und Schulden (+1 Bürde):
\begin{itemize}[leftmargin=*]
\item \textit{Schuld gegenüber einem mächtigen Patron, der Schutz gewährt}
\item \textit{Schuld bei einer Hexe, die dafür gelegentlich Heilung gewährt}
\item \textit{Eid, einen einflussreichen Verbündeten zu unterstützen}
\end{itemize}


\subsection{Schritt 7: Abgeleitete Werte}\label{subsec:schritt-7:-abgeleitete-werte}

Berechne die abgeleiteten Werte für den Charakter und trage sie auf dem Charakterbogen ein.

\subsubsection{Wundenschwelle}
$3 + \text{Körper}$\index{Wundenschwelle}

Die Anzahl an Wunden, die ein Charakter erleiden kann, bevor er den Zustand ``Sterbend'' erhält.

\subsubsection{Bürdenschwelle}
$5 + \lfloor \text{Wille} / 2 \rfloor$

Die maximale Bürde, die ein Charakter tragen kann, bevor er einen \textit{Wendepunkt} erreicht.

\subsubsection{Initiative}
$30 + (\text{Geschick} \times 10) + (\text{Wahrnehmung} \times 10)$\index{Initiative}

Bestimmt die Reihenfolge im Konflikt und bei zeitkritischen Situationen.

\subsubsection{Omen}
$2 + \lfloor \text{Glaubensniveau} / 2 \rfloor$

Die maximale Anzahl an Omen-Punkten, die der Gruppe zur Verfügung stehen.
Diese werden nicht von allen Charakteren summiert, die Gruppe hat gemeinsam diese Anzahl an Omen-Punkten.

\subsubsection{Zauberverstärkung}
$\text{Magieniveau}$\index{Zauberverstärkung}

Der Bonus, der auf magische Effekte angewendet wird.

\subsubsection{Anrufungswert}
$\text{Glaubensniveau} + \text{Ritus}$\index{Anrufungswert}

Der Wert, mit dem göttliche Anrufungen durchgeführt werden.

\subsubsection{Gunstgrenze je Anrufung}
$1 + \lfloor \text{Wille} / 2 \rfloor$

Die maximale Gunst, die für eine einzelne Anrufung ausgegeben werden darf.

\subsection{Schritt 8: Startausrüstung}\label{subsec:schritt-8:-startausrustung}

Wähle aus den vorhandenen Ausrüstungsgegenständen die folgenden Dinge aus und vermerke sie auf deinem Charakterbogen.

\begin{itemize}[leftmargin=*]
\item eine Primärwaffe mit Schaden 1--2
\item optional eine Zweitwaffe mit Schaden 1
\item eine Rüstung mit Schutz 1--2
\item drei Gebrauchsgegenstände
\end{itemize}

\subsection{Schritt 9: Startwerte}\label{subsec:schritt-9:-startwerte}

Die sich verändernden Werte des Charakters haben unterschiedliche Startwerte.
Auf diese werden sie bei der Charaktererschaffung gesetzt.

\begin{itemize}[leftmargin=*]
\item Wunden 0
\item Bürde 0
\item Omen = Maximales Omen
\item Arkana = 3 + Geist
\item Gunst = 3 + Wille
\item Verderbnis = 0
\item Mal\index{Mal}: keins
\end{itemize}

\subsection{Schritt 10: Übernatürlicher Zugang}\label{subsec:schritt-10:-ubernaturlicher-zugang}

Ein Charakter mit Zauberei >= 1 ist magisch begabt.
Je nach Höhe des Zauberei-Werts können Zauberschulen und Zauber gewählt werden.

\textbf{Zauberei:}
\begin{itemize}[leftmargin=*]
\item bei Zauberei >= 1: eine Schule wählen
\item bei Zauberei >= 2: zwei Schulen wählen
\item Startzauber: 2 + Zauberei Rang, maximal 5
\end{itemize}

\textbf{Glaube:}
\begin{itemize}[leftmargin=*]
\item jeder Charakter mit Ritus >= 1 darf Anrufungen sprechen.
\item Anrufungen verbrauchen Gunst und nutzen keine Zauberlisten
\end{itemize}

Ein Charakter darf beide Wege zugleich gehen.
\end{multicols*}
